\documentclass[letterpaper, 12pt]{article}
% ============================================================
%  CONFIGURACION COMUN PARA SOLUCIONARIOS PEP 1
% ============================================================

\usepackage[utf8]{inputenc}
\usepackage[spanish]{babel}
\usepackage{amsmath, amssymb, amsfonts}
\usepackage{geometry}
\usepackage{fancyhdr}
\usepackage{enumitem}
\usepackage{graphicx}
\usepackage{hyperref}
\usepackage{needspace}
\usepackage{parskip}
\usepackage{xcolor}
\usepackage{tcolorbox}
\usepackage{eso-pic}
\usepackage{tikz}
\tcbuselibrary{skins, breakable}

\geometry{letterpaper, margin=1.7cm, top=3cm, bottom=2.5cm, headheight=2cm}

\definecolor{celesteSuave}{HTML}{F0F7FF}
\definecolor{azulFuerte}{HTML}{1A5276}
\definecolor{enlace}{HTML}{1A5276}

\newcommand{\R}{\mathbb{R}}
\newcommand{\Dom}{\operatorname{Dom}}
\newcommand{\Rec}{\operatorname{Rec}}
\newcommand{\Cod}{\operatorname{Cod}}

\newcommand{\logoup}{../../../../../template/images/logo_up.png}
\newcommand{\logousach}{../../../../../template/images/logo_usach.png}
\newcommand{\logofondo}{../../../../../template/images/logo_up.png}
\newcommand{\institucion}{Usach Premium}
\newcommand{\curso}{C\'alculo I}
\newcommand{\autor}{Jaime Riquelme}
\newcommand{\semestre}{1er. Semestre de 2026}
\newcommand{\fraseportada}{``Por y para estudiantes''}
\newcommand{\webinstitucion}{www.usachpremium.cl}
\newcommand{\instagraminstitucion}{https://www.instagram.com/usach.premium}
\newcommand{\transparencia}{0.05}

\hypersetup{
    colorlinks=true,
    linkcolor=enlace,
    urlcolor=enlace,
    citecolor=enlace,
    pdfborder={0 0 0},
}

\pagestyle{fancy}
\fancyhf{}
\fancyhead[L]{\includegraphics[height=1.2cm]{\logoup}}
\fancyhead[R]{%
    \begin{minipage}[b]{0.42\textwidth}
        \raggedleft
        \fontsize{9pt}{11pt}\selectfont
        \textbf{\color{azulFuerte}\institucion} \\
        \curso\ -- Solucionario PEP 1 \\
        \textit{\color{gray}@usach.premium}
    \end{minipage}\hspace{0.1cm}%
    \includegraphics[height=1.2cm]{\logousach}%
}
\fancyfoot[L]{\fontsize{9pt}{11pt}\selectfont \textit{\fraseportada}}
\fancyfoot[R]{\fontsize{9pt}{11pt}\selectfont P\'agina \thepage}
\renewcommand{\headrulewidth}{0.4pt}
\renewcommand{\footrulewidth}{0.4pt}

\fancypagestyle{plain}{
    \fancyhf{}
    \renewcommand{\headrulewidth}{0pt}
    \renewcommand{\footrulewidth}{0pt}
}

\newcommand{\MarcaAgua}{
    \AddToShipoutPicture{
        \AtPageCenter{
            \begin{tikzpicture}[remember picture, overlay]
                \node[opacity=\transparencia] at (0,0)
                    {\includegraphics[width=0.7\textwidth]{\logofondo}};
            \end{tikzpicture}
        }
    }
}

\newtcolorbox{solucionbox}[1]{
    colback=white,
    colframe=azulFuerte,
    coltitle=white,
    colbacktitle=azulFuerte,
    title=\textbf{#1},
    title after break=\textbf{#1} (continuaci\'on),
    fonttitle=\bfseries,
    arc=4pt,
    boxrule=1pt,
    enhanced,
    breakable,
    before={\par\Needspace{0.36\textheight}}
}

\newtcolorbox{respuestabox}{
    colback=celesteSuave,
    colframe=azulFuerte,
    boxrule=0.8pt,
    arc=4pt
}

\newcommand{\portadasolucionario}[2]{
    \thispagestyle{plain}
    \AddToShipoutPicture*{%
        \put(54, 700){\includegraphics[width=2cm]{\logoup}}
        \put(420, 706){\includegraphics[width=5cm]{\logousach}}
    }
    \vspace*{0.5cm}
    \begin{center}
        \color{azulFuerte}
        {\huge \textbf{SOLUCIONARIO}} \\[0.5cm]
        {\Huge \textbf{\curso}} \\[0.3cm]
        {\Large #1} \\[0.3cm]
        {\large #2} \\[1.4cm]
        \begin{tcolorbox}[colback=celesteSuave, colframe=azulFuerte, arc=10pt,
            width=0.82\textwidth, center, boxrule=1.5pt]
            \centering
            \Large\textbf{Desarrollo paso a paso}
        \end{tcolorbox}
    \end{center}
    \vspace{1.4cm}
    \begin{center}
        {\Large \textbf{\semestre}} \\[0.7cm]
        {\large \textbf{\autor\ -- \institucion}}
    \end{center}
    \vfill
    \begin{center}
        {\color{azulFuerte}\hrule height 0.5pt}
        \vspace{0.3cm}
        \begin{minipage}{0.45\textwidth}
            \centering
            \textbf{Instagram:} \\
            \small\href{\instagraminstitucion}{@usach.premium}
        \end{minipage}
        \begin{minipage}{0.45\textwidth}
            \centering
            \textbf{Web:} \\
            \small\href{http://\webinstitucion}{\webinstitucion}
        \end{minipage}
    \end{center}
    \newpage
    \MarcaAgua
}

\newcommand{\respuesta}[1]{
    \begin{respuestabox}
        \textbf{Respuesta:} #1
    \end{respuestabox}
}


\begin{document}
\portadasolucionario{Gu\'ia de Preparaci\'on PEP 1}{Secci\'on III -- Inyectividad, biyectividad e inversa}

\begin{solucionbox}{Ejercicio 3.1}
Sea $f:A\to B$ definida por
\[
    f(x)=(x+1)^2-2.
\]
Determine conjuntos $A$ y $B$ para que $f$ sea biyectiva y $2\in A$. Luego encuentre $f^{-1}$.

\textbf{Soluci\'on:}

La funci\'on corresponde a una par\'abola que abre hacia arriba. Su v\'ertice se encuentra en
\[
    V=(-1,-2).
\]
Para que una funci\'on cuadr\'atica sea inyectiva, debemos restringir su dominio a un lado del v\'ertice.

Como el enunciado pide que $2\in A$, escogemos la rama derecha:
\[
    A=[-1,+\infty).
\]
En este intervalo la funci\'on es creciente, por lo tanto es inyectiva.

Ahora determinamos el recorrido. Si $x\in[-1,+\infty)$, entonces
\[
    x+1\geq 0.
\]
Luego,
\[
    (x+1)^2\geq 0,
\]
y por tanto
\[
    (x+1)^2-2\geq -2.
\]
As\'i,
\[
    \Rec(f)=[-2,+\infty).
\]
Para que $f$ sea sobreyectiva, debemos tomar
\[
    B=\Rec(f)=[-2,+\infty).
\]
Con esto,
\[
    f:[-1,+\infty)\to[-2,+\infty)
\]
es biyectiva.

Para hallar la inversa, escribimos
\[
    y=(x+1)^2-2.
\]
Despejamos $x$:
\[
    y+2=(x+1)^2.
\]
Como $x\in[-1,+\infty)$, se tiene $x+1\geq 0$, por lo que
\[
    \sqrt{y+2}=x+1.
\]
Luego,
\[
    x=-1+\sqrt{y+2}.
\]
Por lo tanto,
\[
    f^{-1}(y)=-1+\sqrt{y+2}.
\]

\respuesta{$A=[-1,+\infty)$, $B=[-2,+\infty)$ y $f^{-1}(y)=-1+\sqrt{y+2}$.}
\end{solucionbox}

\begin{solucionbox}{Ejercicio 3.2}
Sea $k\in\R$ tal que
\[
    f:[-2,1]\to[k,27],
    \qquad
    f(x)=(x-3)^2+2
\]
est\'a bien definida. Demuestre que $f$ es inyectiva, determine $k$ para que sea sobreyectiva y encuentre $f^{-1}$ para ese valor de $k$.

\textbf{Soluci\'on:}

Primero probaremos que $f$ es inyectiva en $[-2,1]$.

Sean $a,b\in[-2,1]$ tales que
\[
    f(a)=f(b).
\]
Entonces
\[
    (a-3)^2+2=(b-3)^2+2.
\]
Restando $2$:
\[
    (a-3)^2=(b-3)^2.
\]
Luego,
\[
    |a-3|=|b-3|.
\]
Como $a,b\in[-2,1]$, se tiene
\[
    a-3<0,\qquad b-3<0.
\]
Por lo tanto,
\[
    |a-3|=-(a-3),\qquad |b-3|=-(b-3).
\]
As\'i,
\[
    -(a-3)=-(b-3),
\]
lo que implica
\[
    a-3=b-3
    \quad\Longrightarrow\quad
    a=b.
\]
Por lo tanto, $f$ es inyectiva.

Ahora hallamos el recorrido. Como $f$ es decreciente en $[-2,1]$, calculamos los extremos:
\[
    f(-2)=(-2-3)^2+2=25+2=27,
\]
\[
    f(1)=(1-3)^2+2=4+2=6.
\]
Por lo tanto,
\[
    \Rec(f)=[6,27].
\]
Para que $f$ sea sobreyectiva sobre $[k,27]$, se debe cumplir
\[
    [k,27]=[6,27],
\]
es decir,
\[
    k=6.
\]

Para hallar la inversa, consideramos
\[
    y=(x-3)^2+2,
    \qquad x\in[-2,1].
\]
Despejamos:
\[
    y-2=(x-3)^2.
\]
Como $x\in[-2,1]$, entonces $x-3<0$, por lo que
\[
    \sqrt{y-2}=|x-3|=-(x-3)=3-x.
\]
Luego,
\[
    x=3-\sqrt{y-2}.
\]

\respuesta{$k=6$ y $f^{-1}(y)=3-\sqrt{y-2}$, con $f^{-1}:[6,27]\to[-2,1]$.}
\end{solucionbox}

\begin{solucionbox}{Ejercicio 3.3}
Considere
\[
    q(x)=2-\sqrt{9-(x-4)^2}.
\]
Determine $\Dom(q)$. Luego escoja un subconjunto $A_1\subseteq \Dom(q)$ donde $q$ sea inyectiva y encuentre una funci\'on inversa indicando dominio y recorrido.

\textbf{Soluci\'on:}

Para hallar el dominio, pedimos que el radicando sea mayor o igual que cero:
\[
    9-(x-4)^2\geq 0.
\]
Esto equivale a
\[
    (x-4)^2\leq 9.
\]
Aplicando ra\'iz:
\[
    |x-4|\leq 3.
\]
Por lo tanto,
\[
    -3\leq x-4\leq 3.
\]
Sumando $4$:
\[
    1\leq x\leq 7.
\]
Luego,
\[
    \Dom(q)=[1,7].
\]

La expresi\'on es sim\'etrica respecto de $x=4$, por lo que no es inyectiva en todo $[1,7]$. Escogemos una de sus ramas:
\[
    A_1=[4,7].
\]
En este intervalo, al aumentar $x$, el valor $(x-4)^2$ aumenta, por lo que $9-(x-4)^2$ disminuye, su ra\'iz disminuye y entonces $q(x)=2-\sqrt{9-(x-4)^2}$ aumenta. Por lo tanto, $q$ es inyectiva en $[4,7]$.

Calculamos el recorrido en esta rama:
\[
    q(4)=2-\sqrt{9}=2-3=-1,
\]
\[
    q(7)=2-\sqrt{9-9}=2.
\]
As\'i,
\[
    \Rec(q|_{[4,7]})=[-1,2].
\]

Ahora despejamos la inversa. Sea
\[
    y=2-\sqrt{9-(x-4)^2}.
\]
Entonces
\[
    \sqrt{9-(x-4)^2}=2-y.
\]
Elevando al cuadrado:
\[
    9-(x-4)^2=(2-y)^2.
\]
Luego,
\[
    (x-4)^2=9-(2-y)^2.
\]
Como $x\in[4,7]$, se tiene $x-4\geq 0$, por lo tanto
\[
    x-4=\sqrt{9-(2-y)^2}.
\]
As\'i,
\[
    x=4+\sqrt{9-(y-2)^2}.
\]

\respuesta{$\Dom(q)=[1,7]$. Tomando $A_1=[4,7]$, se tiene $q^{-1}(y)=4+\sqrt{9-(y-2)^2}$, con $q^{-1}:[-1,2]\to[4,7]$.}
\end{solucionbox}

\begin{solucionbox}{Ejercicio 3.4}
Considere
\[
    F:]-2,2[\to\R,
    \qquad
    F(x)=\frac{x}{2-|x|}.
\]
Pruebe que $F$ es biyectiva y determine $F^{-1}$.

\textbf{Soluci\'on:}

Estudiamos la funci\'on por tramos.

Si $x\in[0,2)$, entonces $|x|=x$, y
\[
    F(x)=\frac{x}{2-x}.
\]
Esta funci\'on toma valores en $[0,+\infty)$. Para verlo, si $y=\frac{x}{2-x}$, entonces
\[
    y(2-x)=x,
\]
\[
    2y=xy+x=x(y+1),
\]
por lo tanto
\[
    x=\frac{2y}{1+y},\qquad y\geq 0.
\]

Si $x\in(-2,0)$, entonces $|x|=-x$, y
\[
    F(x)=\frac{x}{2+x}.
\]
Esta funci\'on toma valores en $(-\infty,0)$. Si
\[
    y=\frac{x}{2+x},
\]
entonces
\[
    y(2+x)=x,
\]
\[
    2y+xy=x,
\]
\[
    2y=x(1-y),
\]
por lo tanto
\[
    x=\frac{2y}{1-y},\qquad y<0.
\]

Como el primer tramo cubre $[0,+\infty)$ y el segundo cubre $(-\infty,0)$, el recorrido total es $\R$. Adem\'as, cada tramo es inyectivo y sus recorridos no se intersectan. Por lo tanto, $F$ es biyectiva.

Unificando las dos expresiones:
\[
    F^{-1}(y)=
    \begin{cases}
        \dfrac{2y}{1-y}, & y<0,\\[1ex]
        \dfrac{2y}{1+y}, & y\geq 0.
    \end{cases}
\]
Como
\[
    1+|y|=
    \begin{cases}
        1-y, & y<0,\\
        1+y, & y\geq 0,
    \end{cases}
\]
podemos escribir
\[
    F^{-1}(y)=\frac{2y}{1+|y|}.
\]

\respuesta{$F$ es biyectiva y $F^{-1}(y)=\dfrac{2y}{1+|y|}$, con $F^{-1}:\R\to]-2,2[$.}
\end{solucionbox}

\begin{solucionbox}{Ejercicio 3.5}
Sea
\[
    r:\R-\{-2\}\to\R-\{3\},
    \qquad
    r(x)=\frac{3x-1}{x+2}.
\]
Justifique que $r$ es biyectiva y encuentre $r^{-1}$.

\textbf{Soluci\'on:}

Para hallar la inversa, escribimos
\[
    y=\frac{3x-1}{x+2}.
\]
Multiplicamos por $x+2$:
\[
    y(x+2)=3x-1.
\]
Desarrollando:
\[
    xy+2y=3x-1.
\]
Agrupamos los t\'erminos con $x$:
\[
    xy-3x=-1-2y.
\]
Factorizamos:
\[
    x(y-3)=-(1+2y).
\]
Como el codominio excluye $y=3$, podemos dividir por $y-3$:
\[
    x=\frac{-(1+2y)}{y-3}
     =\frac{1+2y}{3-y}.
\]
Entonces
\[
    r^{-1}(y)=\frac{1+2y}{3-y}.
\]

Esta expresi\'on est\'a definida para todo $y\neq 3$, y toma valores distintos de $-2$. En efecto, si
\[
    \frac{1+2y}{3-y}=-2,
\]
entonces
\[
    1+2y=-6+2y,
\]
lo cual es imposible. As\'i, la inversa va desde $\R-\{3\}$ hacia $\R-\{-2\}$.

Como existe una funci\'on inversa definida en todo el codominio, $r$ es biyectiva.

\respuesta{$r$ es biyectiva y $r^{-1}(y)=\dfrac{1+2y}{3-y}$, con $r^{-1}:\R-\{3\}\to\R-\{-2\}$.}
\end{solucionbox}

\end{document}
